\section{Introduction}
\label{sec:intro}

\subsection{Motivation}

IP geolocation underpins a broad set of network operator decisions: ISPs steer traffic toward
nearby peers, CDNs select edge servers, regulators enforce jurisdictional content rules, and
capacity planners forecast regional demand. At Internet scale, locating millions of IP addresses
on a recurring basis is unavoidable. Commercial geolocation databases are black-box, expensive,
and raise privacy concerns. Operators need a measurement-light, in-house alternative they can
actually trust.

The dominant measurement-based technique is Constraint-Based Geolocation (CBG)~\cite{gueye2006constraint,katz2006towards}: collect round-trip
times (RTTs) from a set of \emph{vantage points} (VPs) whose locations are known, convert each
RTT to a distance bound, and intersect the resulting geometric constraints to estimate the
target's position. CBG is entirely passive-capable, requires no cooperation from the target
organization, and scales to millions of addresses. However, since Katz-Bassett et al.~\cite{katz2006towards}
and Laki et al.~\cite{laki2011spotter}, the literature has produced several
incompatible algorithmic variants---Vanilla, Million-scale~\cite{poese2011ip}, Octant~\cite{wong2007octant},
Spotter~\cite{laki2011spotter}---without a controlled, operator-grounded comparison of when
and why each one works.

\subsection{The Operator Setting Is Not the General Problem}
\label{sec:intro-operator}

The operator setting differs from general-purpose IP geolocation in three structural ways that
prior work conflates:

\textbf{Free vantage points.} Operators already collect passive RTTs at core infrastructure
across many physical sites---a no-cost fleet of probes whose locations are ground truth.

\textbf{Bounded answer space.} For hypergiants, CDNs, and peering partners, operators already
know the set of possible locations, because that is where they built the interconnection. An
Akamai IP observed at a mobile operator's core does not live at an arbitrary latitude/longitude;
it lives at one of $N$ known data-center sites.

\textbf{Evaluation as classification, not regression.} Because the answer space is finite,
success is no longer absolute coordinate error---it is \emph{classification accuracy} over the
known candidate set. A CBG prediction is correct if and only if the emitted coordinate snaps to
the right site when projected onto the candidate set. This is easier than continuous regression
(finite, tolerant of bounded error) and more useful (matches the operator's actual decision).

We operationalize this by clustering ground-truth target coordinates into answer cells of radius
$R = 50$\,km via complete-linkage agglomerative clustering. A prediction is correct iff it snaps
to the truth's centroid. This yields 257 globally distinct cells over 713 RIPE Atlas anchors,
32 cells for US-targeted evaluation, and 21 for Germany-targeted evaluation.

\subsection{Why Existing Methods Are Not Enough}
\label{sec:intro-challenges}

A natural reaction: \emph{``If operators know their peers' possible sites, why not use Geofeed
or rDNS as the canonical signal?''} In practice those signals leave a large, structurally biased
residual:

\begin{itemize}[leftmargin=*]
  \item \textbf{Geofeed ingestion is fragmented.} No single ingestion endpoint exists. Building a
    pipeline that \emph{finds} feeds across WHOIS remarks and self-hosted URLs is itself substantial
    engineering.
  \item \textbf{Staleness.} Operators do not republish on every relocation or capacity expansion,
    so even reachable data can be wrong.
  \item \textbf{Manipulation.} Both signals are self-reported and can be purposefully misreported
    for competitive advantage; they cannot be trusted without external validation.
  \item \textbf{Coverage gaps.} For entire ASNs---including ones whose traffic operators care
    about---\emph{no} rDNS or Geofeed signal is available at all.
\end{itemize}

The coverage gap is decisive: when canonical signals are missing, the operator still needs a
location estimate. Latency measurement is the only remaining lever that runs entirely on the
operator's own infrastructure and requires no cooperation from the target.

Within that lever we separate two distinct uses. The \textbf{physics-validation primitive} uses
the speed-of-Internet bound ($\approx \tfrac{2}{3}c$) to validate claimed locations: a claim is
\emph{physics-consistent} iff it lies inside the speed-of-Internet circle of every VP. This is
calibration-free and bootstraps labeling from Geofeed/rDNS records. The \textbf{calibrated CBG
estimators} (Vanilla, Octant, Spotter) fit per-VP latency-to-distance models from labeled
RTT-distance pairs and emit point estimates for IPs without canonical records.

This separation yields two operator workflows. In the \textbf{Anchored track}, any rDNS or
Geofeed coverage exists for the ASN: speed-of-Internet validates the claims; the validated subset
becomes training labels; calibrated CBG variants train per-ASN on those labels; the operator gets
verification for covered IPs and prediction for the uncovered residual. In the
\textbf{Unanchored track}, zero canonical coverage is available: only the speed-of-Internet
(Million-scale) prediction is possible, evaluated against the bounded answer space.

\subsection{Thesis and Scope}
\label{sec:intro-thesis}

We present the first cross-variant CBG benchmark from the network-operator perspective, run on
two datasets:

\begin{itemize}[leftmargin=*]
  \item A \textbf{public RIPE-based dataset} with RTTs from three ASN-defined fleets---Amazon
    AWS~(AS16509, 30\,VPs globally), AT\&T~(AS7018, 125\,US VPs), Vodafone~DE~(AS3209,
    164\,EU VPs)---against 713 RIPE Atlas anchors. This dataset is reproducible and serves as
    the public stress test.
  \item A \textbf{proprietary operator dataset} with one operator's user-plane/mobile-core RTTs
    toward a single target ASN across 20+ US clusters. This is the clean operator regime: RTTs
    from one operator ASN to one target ASN share the same peering, backbone, and interconnection
    topology that CBG's calibration assumes.
\end{itemize}

We decompose each CBG variant into a three-phase pipeline (Latency-to-Distance $\to$
Multilateration $\to$ Centroid), evaluate all phases under both continuous and discrete criteria,
identify when and why each variant succeeds or fails, characterize production costs, and release
an open-source implementation.

\subsection{Contributions}
\label{sec:intro-contributions}

\begin{enumerate}[leftmargin=*]
  \item \textbf{Operator-grounded benchmark.} First controlled cross-variant CBG evaluation
    under operator-realistic VP setups (ASN-scoped fleets, matched VP topology between datasets)
    and an operator-relevant evaluation metric (bounded-candidate classification at $R = 50$\,km).

  \item \textbf{Unified composable framework.} A three-phase abstraction spanning
    Vanilla/Million-scale/Octant/Spotter, implemented with interchangeable phase components
    and deployed under a two-track taxonomy (Anchored / Unanchored).

  \item \textbf{Tolerance dividend.} We quantify the gap between coordinate-error and
    classification accuracy as the \emph{tolerance dividend}: in the matched-regional US setup
    (AS7018, 35\,km median VP distance), 53--68\% of all correct answers are won only by
    answer-space tolerance, invisible to any coordinate-error metric. The dividend scales with
    VP-to-target distance and shrinks to 11--33\% in the DE setup (AS3209, 4.6\,km median).

  \item \textbf{Failure taxonomy and EXCLUSIVE-but-correct phenomenon.}
    A two-layer geometric taxonomy---EMPTY\_REGION, EXCLUSIVE\_REGION, INCLUSIVE
    misclassification---partitions every wrong prediction by \emph{where} in the three-phase
    pipeline the error originates. Contrary to intuition, EXCLUSIVE is \emph{not} zero for
    annulus-based variants in proximity-sufficient setups: Octant exhibits $\approx$50\%
    EXCLUSIVE (joint annulus intersection drifts beyond D\_truth's disk) and Spotter $\approx$89--99\%
    EXCLUSIVE. Yet classification accuracy survives: the centroid of an EXCLUSIVE region can still
    snap to the correct answer cell, because Voronoi cells extend beyond D\_truth's 50\,km disk.
    We show that Octant's entire tolerance dividend ($\approx$27\,pp in the US setup) is carried by
    these EXCLUSIVE-but-still-correct predictions. Disk-based variants (Vanilla, Million-scale)
    remain near-zero EXCLUSIVE (0--14\%), with failures dominated by INCLUSIVE misclassification.

  \item \textbf{Per-ASN calibration evidence.} Empirical evidence that pooled-ASN training
    degrades accuracy; per-ASN calibration is the correct methodological default for operator
    deployments.

  \item \textbf{Production-cost characterization.} The centroid phase dominates cost by two
    orders of magnitude: Monte Carlo medoid takes 190--390\,ms per target; geometric centroid
    takes $\approx 0.25$\,ms. Replacing Monte Carlo with geometric centroid in Octant yields
    $\approx 20\times$ throughput with equal or better accuracy.
\end{enumerate}
